% Generated draft: requires mathematical and source review.
\documentclass[11pt]{article}
\usepackage[margin=25mm]{geometry}
\usepackage{fontspec}
\setmainfont{Latin Modern Roman}
\usepackage{amsmath,amssymb,mathtools,unicode-math}
\setmathfont{Latin Modern Math}
\usepackage{enumitem,needspace,xurl,hyperref}
\setlength{\parindent}{0pt}
\setlength{\parskip}{5pt}
\setlength{\emergencystretch}{3em}
\newcommand{\N}{\mathbb{N}}
\newcommand{\Z}{\mathbb{Z}}
\newcommand{\Q}{\mathbb{Q}}
\newcommand{\R}{\mathbb{R}}
\newcommand{\C}{\mathbb{C}}
\newcommand{\F}{\mathbb{F}}
\newcommand{\A}{\mathbb{A}}
\newcommand{\PP}{\mathbb{P}}
\newcommand{\E}{\mathbb{E}}
\newcommand{\eps}{\varepsilon}
\newcommand{\dd}{\,\mathrm{d}}
\begin{document}
\section*{UnsolvedMath ID: 30005897}
\textbf{Source problem code:} catalogue entry
\subsection{Definitions and mathematical statement}
The supplied record asks whether a shadowing characterization for dissipative composition operators extends when bounded distortion is no longer assumed. It also asks specifically whether generalized hyperbolicity and shadowing remain equivalent. The precise underlying function space, measure-theoretic hypotheses, and characterization referred to as [1] are not reproduced in the record; consequently, the complete original formulation cannot be recovered from the supplied text alone.

The following definitions give a conditional formulation under the usual invertible Banach-space conventions. They do not establish that these are exactly the conventions of the cited report. Let \((X,\Sigma,\mu)\) be a measure space, let \(E\) be a Banach space of scalar-valued measurable functions modulo equality almost everywhere, and let \(f\colon X\to X\) induce a well-defined bounded invertible composition operator
\[
T=C_f\colon E\to E,\qquad C_f\varphi=\varphi\circ f.
\]
The record does not specify \(E\), its scalar field, or the hypotheses on \(f\) ensuring these properties. In a common invertible dissipative setting, there is a measurable wandering set \(W\) such that the sets \(f^n(W)\), \(n\in\Z\), are pairwise disjoint modulo null sets and
\[
X=\bigcup_{n\in\Z}f^n(W)\quad\text{modulo null sets}.
\]
Any additional requirements on \(W\), such as finiteness or positivity of its measure, must be recovered from the source.

For a bounded invertible operator \(T\) on a Banach space \(E\), call \(T\) generalized hyperbolic if there are closed subspaces \(M,N\subseteq E\), constants \(K\geq1\) and \(0<r<1\), such that
\[
E=M\oplus N,\qquad T(M)\subseteq M,\qquad T^{-1}(N)\subseteq N,
\]
and, for every integer \(n\geq0\),
\[
\|T^n u\|\leq Kr^n\|u\|\quad(u\in M),\qquad
\|T^{-n}v\|\leq Kr^n\|v\|\quad(v\in N).
\]
Under the two-sided convention, \(T\) has the shadowing property if for every \(\varepsilon>0\) there exists \(\delta>0\) such that every sequence \((x_n)_{n\in\Z}\subseteq E\) satisfying
\[
\|Tx_n-x_{n+1}\|<\delta\qquad(n\in\Z)
\]
admits \(x\in E\) with
\[
\|T^n x-x_n\|<\varepsilon\qquad(n\in\Z).
\]
The record does not specify whether this two-sided convention is intended.

With the original ambient hypotheses and conventions fixed, let \(\mathcal D\) denote the intended class of dissipative composition operators after removal of the bounded-distortion hypothesis. The requested equivalence is
\[
\text{For every }C_f\in\mathcal D,\qquad
C_f\text{ is generalized hyperbolic}
\quad\Longleftrightarrow\quad
C_f\text{ has the shadowing property}.
\]
Here removal of the hypothesis permits operators without bounded distortion; it need not mean restricting to operators that fail bounded distortion. The source wording leaves this distinction unresolved. The additional question asks whether the full characterization of shadowing cited as [1] remains valid in this enlarged setting. Since that characterization and the identification of [1] are not explicitly supplied, this additional question cannot be stated as an explicit mathematical equivalence from the record alone.
\subsection{Short English statement}
For the dissipative composition operators intended in the source, does removing the bounded-distortion assumption preserve the equivalence between generalized hyperbolicity and shadowing? The original passage also asks whether an earlier full characterization of shadowing extends. The supplied record omits the exact ambient hypotheses and that characterization, so the displayed formulation is conditional on recovering them.
\subsection{Sources}
\begin{itemize}
\item Local UnsolvedMath registry, permanent ID \(30005897\), title \emph{Shadowing Without Bounded Distortion}, source problem code \texttt{catalogue entry}. The supplied record belongs to the Oberwolfach Reports Open Problems collection and assigns the problem to 2024; no individual proposer is supplied.
\item \emph{Mini-Workshop: New Horizons in Linear Dynamics, Universality, and the Invariant Subspace Problem} (2024), Oberwolfach Reports, pp.\ 1069--1096, DOI \url{https://doi.org/10.4171/owr/2024/19}. The supplied original excerpt locates the open-problem discussion on p.\ 1081. This report is the supplied primary-source attribution; it has not been independently consulted here.
\item Supplied literature reference: D'Aniello, Darji and Maiuriello, \url{https://doi.org/10.1016/j.jde.2021.06.038}. The record does not supply the article title, theorem statement, or an explicit identification of this reference with [1] in the original excerpt.
\item Supplied literature reference: Oberwolfach Report 19/2024, \url{https://doi.org/10.4171/owr/2024/19}.
\item Literature assessment dated exactly \texttt{2026-08-22}: the supplied curation labels the problem open and reports that the equivalence and full shadowing characterization are proved under bounded distortion, while the setting without that assumption remains unresolved. This is supplied dated curation, not an independently verified literature assessment.
\item The registry labels the statement \texttt{corrected\_verified} and reports reconstruction and checking against the official Oberwolfach report. These are supplied provenance assertions, not independent verification of the formulation or its mathematical status.
\end{itemize}
\end{document}
